\section{Motivation}
Zur Steigerung der Wirtschaftlichkeit landwirtschaftlicher Betriebe sind die Personalkosten ein entscheidender Faktor. Im Getreideanbau werden Zug- und Anbaugeräte seit Jahrzehnten immer größer, um die gleiche Fläche mit weniger Personal bearbeiten zu können. Der Gemüsebau ist jedoch weiterhin auf viel menschliche Handarbeit angewiesen. Dafür werden in Deutschland Saisonarbeitskräfte aus dem Ausland angeworben. Dabei bereiten sowohl der Arbeitskräftemangel als auch deren Mindestlohn den Landwirten Schwierigkeiten \cite{arbeitskräfte}. In den letzten Jahren hat es viele Verbesserungen im Bereich des maschinellen Lernens gegeben. Durch Leistungssteigerungen in der Hardware konnten bei Computer Vision, Deep Learning Frameworks wie YOLO oder RFDETR, aber auch bei Imitation und Reinforcement Learning große Erfolge erzielt werden. Deshalb hält es die Firma AI.Land für absehbar, dass Roboterarme und -hände in naher Zukunft einfache, sich wiederholende Tätigkeiten von Menschen übernehmen können. Diese können auf selbstfahrenden Plattformen wie dem etaROB oder DAVEGI mobile von AI.Land montiert werden, sodass sie vollautomatisiert unterschiedliche Aufgaben übernehmen können. Denkbar wäre das Ernten von Gemüse, das Entfernen von Unkraut auch dicht an den Nutzpflanzen und das Ausbringen von Herbiziden und Pestiziden nur dort, wo sie wirklich nötig sind. Zu einem autonomen Roboter gehört auch ein möglichst autarkes Energiesystem. Photovoltaik (PV) kann lokal, aber unregelmäßig Energie bereitstellen. Aus dem PV-Strom kann Wasserstoff mittels Elektrolyseuren ebenfalls lokal hergestellt werden und dem Landwirt somit Energieautarkie für einen Feldroboter liefern. \\ 

Das Ziel dieser Bachelorarbeit ist es, die verschiedenen Komponenten eines Energiesystems für den Feldroboter DAVEGI mobile mithilfe einer Computersimulation auszulegen. Betrachtet werden eine PEM-Brennstoffzelle mit einem Wasserstofftank variabler Größe, eine variable Kapazität an $LiFePO_4$-Akkumulatoren und eine variable Fläche an PV-Modulen. Die Auslegung beinhaltet die Auswahl relevanter Systemparameter, das Festlegen einer Regelungsstrategie und eine Pareto-Front-basierte Auswahl sinnvoller Kombinationsmöglichkeiten der variablen Größen. Die zu optimierenden Größen sind die Auslastung des Feldroboters bei möglichst geringem Gewicht. Dabei werden unterschiedliche Monate im Jahr sowie unterschiedliche Tankintervalle betrachtet. Eine möglichst hohe Auslastung ist für die Amortisation des Feldroboters erstrebenswert. Die Tankintervalle sollen möglichst groß sein, um die Autonomie des Feldroboters zu maximieren, da ein häufiges Betanken den Arbeitsaufwand der Landwirte erhöht. Das Gewicht soll zudem so gering wie möglich gehalten werden, um Bodenverdichtung vorzubeugen. Bodenverdichtung reduziert das Bewurzelungsvermögen vieler Pflanzen im Ackerboden, die Belüftung des Bodens, die Bodenaktivität durch Kleinstlebewesen und die schnelle Wasseraufnahme bei Starkregen, was zu Erosion führen kann. \\

Kommentar (streiche ich wieder): Dampfreformation aus Biogas ist im Zuge der Energiewende mehr schädlich als sinnvoll. weshalb ich es nicht erwähne. \\





